%%
%% fall-lecture20.tex
%% 
%% Made by Alex Nelson <pqnelson@gmail.com>
%% Login   <alex@lisp>
%% 
%% Started on  2026-11-12T09:13:37-0800
%% Last update 2026-11-12T09:13:37-0800
%% 

\lecture{}

\begin{recall}
We are studying ``quasi 1-dimensional flows'' (which is especially
relevant in aerodynamics).
We introduced, for many physical quantities $X$ associated
relationships with the Mach number like $(1/X)\D X\sim(1\Mach)\D\Mach$,
which we integrated to obtain a ratio between the quantity $X$ and its
``stagnation quantity'' $X_{0}$ and some function of $X_{0}/X\sim(1+\beta\Mach^{2})^{C}$:
\begin{equation}
\left.\begin{array}{r}
\displaystyle{\frac{\D u}{u} = \frac{1}{1+\beta\Mach^{2}}\frac{\D\Mach}{\Mach}}\\[12pt]
\displaystyle{\frac{\D a}{a} = \frac{-\beta\Mach^{2}}{1+\beta\Mach^{2}}\frac{\D\Mach}{\Mach}}\\[12pt]
\displaystyle{\frac{\D\rho}{\rho}=\frac{-\Mach^{2}}{1+\beta\Mach^{2}}\frac{\D\Mach}{\Mach}}\\[12pt]
\displaystyle{\frac{\D h}{2(h-h_{0})}=\frac{1}{1+\beta\Mach^{2}}\frac{\D\Mach}{\Mach}}\\[12pt]
\dots
\end{array}\right\}\xrightarrow{\Gamma=\text{const.}}
\left\{\begin{array}{l}
\displaystyle{\frac{\rho_{0}}{\rho}=(1+\beta\Mach^{2})^{1/(2\beta)}}\\[12pt]
\displaystyle{\frac{a_{0}}{a}=(1+\beta\Mach^{2})^{1/2}}\\[12pt]
\displaystyle{\frac{h_{0}}{h}=(1+\beta\Mach^{2})}\\[12pt]
\displaystyle{\frac{p_{0}}{p}=(1+\beta\Mach^{2})^{(2\beta+1)/(2\beta)}}\\[12pt]
\dots
\end{array}\right.
\end{equation}
This was done when $\Gamma$ is constant (equivalently, $\beta$
constant). The $\rho_{0}$, $a_{0}$, $p_{0}$, $h_{0}$ are stagnation
quantities obtained via following isentroopes to the state where the
flow is not moving. We also have the \define{Sonic Area} $A^{*}$
satisfying
\begin{equation}
\frac{A}{A^{*}}=\frac{1}{\Mach^{2}}\left(\frac{1+\beta\Mach^{2}}{1+\beta}\right)^{(1+\beta)/\beta}.
\end{equation}
Recall, we defined
\begin{equation}
\beta:=\frac{\rho}{a}\left(\frac{\partial a}{\partial\rho}\right)_{S},
\end{equation}
so $\Gamma=1+\beta$.
\end{recall}

\begin{assume}
We have a perfect gas (so $c_{V}$ is constant, and $c_{p}$ is constant).
\end{assume}

\begin{node}[Calculating $\beta$ for perfect gases]
Recall Eq~\eqref{eq:me101a:fall2026:lec18:pressure-ratio-propto-density-ratio},
\begin{equation}
\left(\frac{p}{p_{0}}\right)=\left(\frac{\rho}{\rho_{0}}\right)^{\gamma},
\end{equation}
which means
\begin{subequations}
  \begin{align}
a^{2} &= \left(\frac{\partial p}{\partial\rho}\right)_{S}\quad\mbox{by def}\\
&= \frac{p_{0}}{\rho_{0}^{\gamma}}\gamma\rho^{\gamma-1},
  \end{align}
\end{subequations}
then we can compute
\begin{subequations}
  \begin{align}
\beta&=\frac{\rho}{a}\left(\frac{\partial a}{\partial\rho}\right)_{S}\quad\mbox{by def}\\
&=\frac{\rho}{2a^{2}}\frac{p_{0}}{\rho_{0}^{\gamma}}\gamma(\gamma-1)\rho^{\gamma-2}\\
&=\frac{1}{2a^{2}}\frac{p_{0}}{\rho_{0}^{\gamma}}\gamma(\gamma-1)\rho^{\gamma-1},
  \end{align}
\end{subequations}
which gives us (combining these two calculations),
\begin{equation}
\boxed{\beta=\frac{\gamma-1}{2}\quad\mbox{for perfect gases}.}
\end{equation}
This should make sense since $\Gamma=(\gamma+1)/2$ for perfect gases.
\end{node}

\begin{node}[Stagnation quantities for perfect gases]
Using $\beta=(\gamma-1)/2$ for perfect gases, we can rewrite the
stagnation quantity ratios for perfect gases as:
\begin{subequations}
  \begin{align}
\frac{\rho_{0}}{\rho} &= \left(1+\frac{\gamma-1}{2}\Mach^{2}\right)^{1/(\gamma-1)}\\
\frac{a_{0}}{a} &= \left(1+\frac{\gamma-1}{2}\Mach^{2}\right)^{1/2}\\
\frac{h_{0}}{h} &= \left(1+\frac{\gamma-1}{2}\Mach^{2}\right)\\
\frac{p_{0}}{p} &= \left(1+\frac{\gamma-1}{2}\Mach^{2}\right)^{\gamma/(\gamma-1)}\\
\intertext{and,}
\frac{A}{A^{*}} &= \frac{1}{\Mach}\left[\frac{2}{\gamma+1}\left(1+\frac{\gamma-1}{2}\Mach^{2}\right)\right]^{(\gamma+1)/[2(\gamma-1)]}.
  \end{align}
\end{subequations}
\end{node}

\begin{observation}
Everything has been expressed as a function of the Mach number (and
the constant\footnote{It's a perfect gas! Of course
$\gamma=c_{p}/c_{V}$ is constant!} $\gamma$). 
\end{observation}

\begin{node}
Suppose we have a known stagnation condition $\rho_{0}$, $p_{0}$,
$T_{0}$, \dots. As $\Mach\to\infty$ (we speed things up), we see
\begin{subequations}
\begin{equation}
h,\rho,a,p\xrightarrow{\Mach\to\infty} 0,
\end{equation}
\begin{equation}
\frac{A}{A^{*}}\xrightarrow{\Mach\to\infty}\infty,
\end{equation}
and
\begin{equation}
u\xrightarrow{\Mach\to\infty}\sqrt{\frac{2}{\gamma-1}}a_{0}.
\end{equation}
\end{subequations}
That last result concerning $u$ seems surprising, let us check:
\begin{subequations}
  \begin{align}
\frac{u}{a_{0}} &= \frac{u}{a}\frac{a}{a_{0}}\\
&=\Mach\frac{a}{a_{0}}\quad\mbox{by def of Mach number}\\
&=\Mach\left(1+\frac{\gamma-1}{2}\Mach^{2}\right)^{-1/2}\\
&\xrightarrow{\Mach\to\infty}\sqrt{\frac{2}{\gamma-1}}.
  \end{align}
\end{subequations}
Coincidentally, this is the $a\to0$ limit.

\textsc{Moral}: There is a limit to how
fast our fluid can move if we speed it up isentropically. (If there
was no limit, then $u\to\infty$ as $\Mach\to\infty$.)
\end{node}

\begin{example}
Air has $T_{0}=15^{\circ}~\mathrm{C}$, $a_{0}=340~\mathrm{m}/\mathrm{s}$,
and so $u_{\text{max}}=760~\mathrm{m}/\mathrm{s}$.
\end{example}

\begin{note}
The perfect gas idealization breaks down at high Mach numbers.
\end{note}

\begin{remark}[About the Sonic condition $\Mach=1$]
For a perfect gas, the sonic speed $u^{*}$ relative to stagnation
speed of sound $a_{0}$ satisfies
\begin{equation}
\frac{u^{*}}{a_{0}}=\sqrt{\frac{2}{\gamma+1}},
\end{equation}
since
\begin{subequations}
  \begin{align}
\frac{u^{*}}{a_{0}} &= \frac{u^{*}}{a^{*}}\frac{a^{*}}{a_{0}}\\
&=\Mach\frac{a^{*}}{a_{0}}\\
&=\frac{a^{*}}{a_{0}}\\
&=\left(1+\beta\Mach^{2}\right)^{-1/2}=\frac{1}{\sqrt{\Gamma}}.
  \end{align}
\end{subequations}
Similarly, we can find the sonic density
\begin{subequations}
  \begin{align}
\frac{\rho_{0}}{\rho^{*}} &= (1+\beta\Mach^{2})^{1/(2\beta)}\\
&= (1+\beta)^{1/(2\beta)}\\
&= \left(1+\frac{\gamma-1}{2}\right)^{1/(\gamma-1)}\\
&= \left(\frac{2}{1+\gamma}\right)^{-1/(\gamma-1)},
  \end{align}
so (inverting this relation, then multiplying both sides by $\rho^{*}$),
\begin{equation}
\boxed{\rho^{*}=\rho_{0}\left(1+\frac{\gamma-1}{2}\right)^{1/(\gamma-1)}.}
\end{equation}
\end{subequations}
Now, we can find the mass flow $\dot{m}=A\rho u$ at sonic speed
condition as
\begin{subequations}
  \begin{align}
\dot{m}^{*} &=A^{*}\rho^{*}u^{*}\\
&= \left(\frac{2}{\gamma+1}\right)^{(\gamma+1)/2(\gamma-1)}a_{0}\rho_{0}A^{*}.
  \end{align}
\end{subequations}
\end{remark}

\begin{aside}
We can see, following Anderson~\cite[\S8.5]{anderson1984fundamentals},
we can answer the question, ``When is a flow compressible?'' We look
at the ratio,
\begin{equation}
\frac{\rho_{0}}{\rho}=\left(1+\frac{\gamma-1}{2}\Mach^{2}\right)^{1/(\gamma-1)}
\end{equation}
then look at the values of $\Mach$ such that this ratio is
``approximately flat'' (i.e., the density does not change more than
5\% or 10\% or smaller), which means that the flow can be treated as
incompressible. For larger values of $\Mach$, this ratio can
compress. We see when $\Mach=0$ that $\rho=\rho_{0}$ so
$\rho_{0}/\rho=1$.
Let us find the $\Mach_{\text{crit}}$ such that
\begin{equation}
\left(1+\frac{\gamma-1}{2}\Mach^{2}_{\text{crit}}\right)^{1/(\gamma-1)}=\frac{19}{20}.
\end{equation}
After some algebra, we find
\begin{equation}
\Mach_{\text{crit}}^{2}=\frac{2}{\gamma-1}\left((20/19)^{\gamma-1}-1\right).
\end{equation}
For air, $\gamma_{\text{air}}\approx1.4$ which gives $\Mach_{\text{crit}}\approx0.329$
which corresponds to velocities $u<109.46~\mathrm{m}/\mathrm{s}$.
To get some sense of things, a Tornado typically has speeds less than
$50~\mathrm{m}/\mathrm{s}$. (Tornadoes are defined as having speeds
greater than $17.78~\mathrm{m}/\mathrm{s}$.)
\end{aside}

\subsection{Normal Shocks}

\begin{note}
Note to self: (1) ``normal shocks'' seem to be a topic appearing in gas
dynamics (at least it appears there in my books on gas dynamics, but
seems absent from my books on hydrodynamics), but it generalizes to
fluid mechanics more broadly.

(2) ``Normal shock'' as opposed to ``oblique shock'' (or ``non-normal shock'').
\end{note}

\begin{node}
We will look at 1-dimensional flow with no area change, so we have a
``parallel tube'' (i.e., a pipe) and a piston at one end. Assume we
have a compressible fluid.
At time $t<0$, the fluid velocity inside the tube is zero $u=0$.

At time $t=0$, the piston starts moving [at constant velocity?].
\begin{center}
\includegraphics{img/img.41}

\includegraphics{img/img.42}

\includegraphics{img/img.43}

\includegraphics{img/img.44}
\end{center}
At positions $x_{1}$, $x_{2}$, $x_{3}$, we can measure $u$ and observe
it steadily increases in a sort of step function:
\begin{center}
  \includegraphics{img/img.45}
\end{center}
This plot also holds for other quantities like $p$, $\rho$, $T$. At
later times, it smooths out (as wave 1 overtakes wave 2, wave 2
overtakes wave 3, and so on), giving us a plot looking more like:
\begin{center}
  \includegraphics{img/img.46}
\end{center}
\end{node}

\begin{node}
Once a shock forms, it is no longer an isentropic process.
The width of the shock $\lambda$ is ``small''.
\end{node}

\begin{node}[Shock fixed frame]
In this example, we can ``follow the shock''. This is a frame of
reference. But in this frame, the shock wave has zero velocity.
\end{node}

\begin{assume}
\begin{enumerate}
\item Steady flow (or, we are in the shock fixed frame)
\item Neglect friction of wall
\item Neglect heat transfer
\end{enumerate}
\end{assume}

\begin{node}
If we write down the control volume integral form for the equations of
motion, we get:
\begin{subequations}
\begin{equation}\label{eq:fall2026:lec20:eq1}
\rho_{2}u_{2}=\rho_{1}u_{1}=\dot{m}
\end{equation}
\begin{equation}\label{eq:fall2026:lec20:eq2}
p_{2}-p_{1}=\rho_{1}u_{1}^{2}-\rho_{2}u_{2}^{2}
\end{equation}
\begin{equation}\label{eq:fall2026:lec20:eq3}
h_{2}-h_{1}=\frac{1}{2}(u_{1}^{2}-u_{2}^{2}).
\end{equation}
\end{subequations}
\end{node}

\begin{node}
The combining Eq~\eqref{eq:fall2026:lec20:eq1} with Eq~\eqref{eq:fall2026:lec20:eq2},
\begin{subequations}\label{eq:fall2026:lec20:eq4}
  \begin{align}
\Delta p &:= p_{2}-p_{1} \\
&=\rho_{1}u_{1}(u_{1}-u_{2})\\
&=-\rho_{1}u_{1}\,\Delta u\\
&=-\dot{m}\,\Delta u.
  \end{align}
\end{subequations}
Then Eq~\eqref{eq:fall2026:lec20:eq1} gives us
\begin{equation}
\frac{\rho_{1}}{\rho_{2}}=\frac{u_{2}}{u_{1}},
\end{equation}
hence
\begin{equation}\label{eq:fall2026:lec20:eq5}
\frac{\Delta u}{u_{1}}=\frac{u_{2}-u_{1}}{u_{1}}=\frac{-\Delta\rho}{\rho_{2}}.
\end{equation}
Plugging Eq~\eqref{eq:fall2026:lec20:eq4} into Eq~\eqref{eq:fall2026:lec20:eq5}
yields
\begin{equation}\label{eq:fall2026:lec20:eq6}
\frac{\Delta p}{\Delta\rho}=u_{1}u_{2}.
\end{equation}
Recall $V=1/\rho$, which gives us
\begin{equation}\label{eq:fall2026:lec20:eq7}
\frac{\Delta p}{\Delta V}=-\dot{m}^{2}
\end{equation}
We claim
\begin{equation}\label{eq:fall2026:lec20:eq8}
\Delta p\,\Delta V=-(\Delta u)^{2}.
\end{equation}
\end{node}

\begin{aside}
As I transcribe these calculations, I do not quite believe it. We see
\begin{subequations}
  \begin{align}
\Delta\rho &= \rho_{2} - \rho_{1}\\
&= \frac{1}{V_{2}}-\frac{1}{V_{1}}\\
&= \frac{V_{1}-V_{2}}{V_{1}V_{2}}\\
&= \frac{1}{V_{1}}\frac{1}{V_{2}}(-1)(V_{2}-V_{1})\\
&= -\rho_{1}\rho_{2}\,\Delta V.
  \end{align}
\end{subequations}
Then we would have
\begin{subequations}
  \begin{align}
\frac{\Delta p}{\Delta\rho} &= \frac{1}{-\rho_{1}\rho_{2}}\frac{\Delta p}{\Delta V}\\
&= u_{1}u_{2}.
  \end{align}
\end{subequations}
Then we see $\dot{m}=\rho_{1}u_{1}=\rho_{2}u_{2}$ gives us
(multiplying both sides by $-\rho_{1}\rho_{2}$):
\begin{equation}
\frac{\Delta p}{\Delta V}=-\rho_{1}\rho_{2}u_{1}u_{2}=-\dot{m}^{2}.
\end{equation}
This is precisely Eq~\eqref{eq:fall2026:lec20:eq7}. Multiplying
through by $\Delta V$ gives us
\begin{equation}
\Delta p=-\dot{m}^{2}\,\Delta V.
\end{equation}
Then we set this equal to Eq~\eqref{eq:fall2026:lec20:eq4}
\begin{subequations}
  \begin{align}
-\dot{m}^{2}\,\Delta V &= \Delta p\\
&=-\dot{m}\,\Delta u\quad\mbox{by Eq~\eqref{eq:fall2026:lec20:eq4}}
  \end{align}
\end{subequations}
which gives us (assuming $\dot{m}\neq0$)
\begin{equation}
\dot{m}\,\Delta V=\Delta u.
\end{equation}
Using this, multiply both sides by $\Delta p$ and using Eq~\eqref{eq:fall2026:lec20:eq4}
\begin{subequations}
  \begin{align}
\dot{m}\,\Delta V\,\Delta p&=\Delta u\,\Delta p\\
&=-\dot{m}(\Delta u)^{2},
  \end{align}
hence dividing through by $\dot{m}$,
  \begin{equation}
\boxed{\Delta V\,\Delta p=-(\Delta u)^{2}}.
  \end{equation}
\end{subequations}
That's precisely Eq~\eqref{eq:fall2026:lec20:eq8}.
\end{aside}

\begin{node}
Compare these quantities to isentropic equations. For example,
Eq~\eqref{eq:fall2026:lec20:eq7}
``corresponds'' analogously to
\begin{equation}
\left(\frac{\partial p}{\partial V}\right)_{S} = -\rho^{2}a^{2}.
\end{equation}
\end{node}

\begin{node}
Now, we have not touched our conservation of Energy Eq~\eqref{eq:fall2026:lec20:eq3}:
\begin{subequations}
  \begin{align}
\Delta h &= \frac{1}{2}(u_{1}^{2}-u_{2}^{2})\\
&= \frac{1}{2}(u_{1}+u_{2})(u_{1}-u_{2})\\
&= \frac{1}{2}(u_{1}+u_{2})(-\Delta u),
  \end{align}
\end{subequations}
then using Eq~\eqref{eq:fall2026:lec20:eq4} and $\dot{m}=\rho_{1}u_{1}$,
\begin{subequations}
  \begin{align}
\Delta h &= \frac{1}{2}(u_{1}+u_{2})\frac{\Delta p}{\rho_{1}u_{1}}\\
&=\frac{1}{2}(\Delta p)\frac{u_{1}+u_{2}}{\rho_{1}u_{1}},
  \end{align}
\end{subequations}
hence
\begin{subequations}
  \begin{align}
\frac{1}{2}(\Delta p)\left(\frac{1}{\rho_{1}}+\frac{1}{\rho_{2}}\right)
&=\frac{1}{2}(\Delta p)(V_{1}+V_{2})\\
&=\frac{1}{2}(V_{1}+V_{2})(\Delta p)\\
&=\overline{V}\,\Delta p
  \end{align}
\end{subequations}
using $\overline{V}=(V_{1}+V_{2})/2$ the average volume.
\end{node}

\begin{node}
We can similarly compute
\begin{subequations}
  \begin{align}
\Delta e &= e_{2}-e_{1}\\
&=\Delta h + \left(\frac{p_{1}}{\rho_{1}}-\frac{p_{2}}{\rho_{2}}\right)\\
&=\dots\\
&=-\overline{p}\,\Delta V,
  \end{align}
\end{subequations}
where $\overline{p}=(p_{1}+p_{2})/2$ is the average pressure.
\end{node}

\begin{remark}
Zucrow and Hoffman~\cite[\S7--3~\textit{et~seq}]{zucrow1976gas}
is a good review for normal shock waves.
\end{remark}